% Paste this file after \begin{document} in your existing Beamer presentation.
% Beamer already loads graphicx. Copy assets/ next to your main .tex file,
% then adjust this path if needed.
\newcommand{\FlowVizKomegaFigure}{assets/komega_snapshots.png}

\begin{frame}{Flow Viz analysis for thermal-source kernels}
  \begin{columns}[T,onlytextwidth]
    \begin{column}{0.64\textwidth}
      \textbf{New analysis workflow}\\[0.5em]
      \begin{itemize}\setlength{\itemsep}{0.65em}
        \item Probe V2 \texttt{.pbin} and compact HDF5 share one named probe-dataset model.
        \item Multi-probe products include histories, FFT/PSD, STFT and envelopes, pulse response, nonlinear coupling, and modal screening.
        \item Typed Asym/Gaus comparisons align time, frequency, physical probe positions, and wavenumber.
        \item Preflight and phase-level logging report inputs, resource use, output artifacts, and failures while a run is active.
      \end{itemize}
    \end{column}
    \begin{column}{0.31\textwidth}
      \textbf{Thermal-source application}\\[0.5em]
      \small
      Compare the energy kernels with the same probes and product definitions. The workflow separates physical differences in source evolution, spectral content, and wave propagation from file-format and plotting choices.\\[1em]
      \footnotesize
      Suggested figure set:\\[0.3em]
      time history, PSD, time-localized f--k map, and matched-time temperature contours.
    \end{column}
  \end{columns}
\end{frame}

\begin{frame}{Time-localized wavenumber diagnostics}
  \begin{columns}[T,onlytextwidth]
    \begin{column}{0.71\textwidth}
      \centering
      \includegraphics[width=\linewidth,height=0.72\textheight,keepaspectratio]{\FlowVizKomegaFigure}
    \end{column}
    \begin{column}{0.25\textwidth}
      \small
      \textbf{What this figure shows}\\[0.5em]
      Each panel resolves the pressure field into frequency and signed wavenumber at a selected pulse stage.\\[0.8em]
      Positive $k$ denotes downstream propagation for $\cos(\omega t-kx)$.\\[0.8em]
      A shared color scale prevents post-pulse noise from appearing artificially strong.\\[1em]
      \footnotesize
      Illustrative verification output, not an Asym/Gaus production result.
    \end{column}
  \end{columns}
\end{frame}
